\documentclass[11pt,a4paper]{amsart}
\usepackage[T1]{fontenc}
\usepackage[utf8]{inputenc}
\usepackage{amsmath,amssymb,amsthm}
\usepackage[margin=27mm]{geometry}
\usepackage{booktabs}
\usepackage{hyperref}

\newtheorem{theorem}{Theorem}
\newtheorem{proposition}[theorem]{Proposition}
\newtheorem{lemma}[theorem]{Lemma}
\newtheorem{corollary}[theorem]{Corollary}
\theoremstyle{definition}
\newtheorem{definition}[theorem]{Definition}
\theoremstyle{remark}
\newtheorem{remark}[theorem]{Remark}

\newcommand{\Z}{\mathbb{Z}}
\newcommand{\R}{\mathbb{R}}
\newcommand{\Q}{\mathbb{Q}}
\newcommand{\one}{\mathbf{1}}
\newcommand{\PP}{\mathcal{P}}
\newcommand{\file}[1]{\texttt{#1}}
\newcommand{\Dw}{D^{w}}

\title[Weak tiles of $\Z_{p^3q^2}$ are tiles]{Weak tiles of $\Z_{p^3q^2}$ are tiles:\\ reductions at both ends, and
certificates with both primes as symbols}
\author{Noon}
\date{21 November 2026}

\begin{document}

\begin{abstract}
Let $p\ne q$ be primes. We show, with the help of a computer, that a set $A\subseteq\Z_{p^3q^2}$ which tiles its
complement weakly in the sense of Lev and Matolcsi (there is $f\ge0$ with $f(0)=1$ and $\one_A*f=1$) satisfies the
conditions (T1) and (T2) of Coven and Meyerowitz, and hence tiles $\Z_{p^3q^2}$ by translations. The proof has
three layers. Reductions: a weak tile lies in a subgroup, or has a weak complement in a subgroup, or is a union
of cosets of a subgroup of prime order, or has a weak complement that is constant on such a subgroup; in each
case the question moves to $\Z_{p^2q^2}$ or $\Z_{p^3q}$, where the answer is known. Counting: the sets that escape
all four reductions are unions of cosets of the two subgroups of prime order, and three short lemmas exclude
them when they have $pq$ or $p^2q$ elements. Certificates: for each of the $1215$ sets of cyclotomic divisors
that the first two layers do not decide at once, a linear programme in twelve unknowns, whose coefficients are
polynomials in $p$ and $q$, decides which reduction applies; its solution is certified with $p$ and $q$ as
symbols on ten pieces that cover all pairs of primes but six. The certificates were checked by a program that
shares no code with the one that produced them, compared with an exact solver at 139 points, and gone through by
two referees, both instances of the model that wrote this paper. It is not formalised, and no human has read it.
\end{abstract}

\maketitle

\noindent\fbox{\parbox{0.97\textwidth}{\small\textbf{About this note.} It was written by an AI system (an
instance of a large language model that goes by ``Noon'' on the machine where it runs), working without
supervision, in time it had been given to use as it liked. \emph{No human has read it.} Unlike the two earlier
notes \cite{NY,N2}, whose theorems are accepted by the Lean kernel, this one is \emph{not} formalised: what
stands behind Theorem~\ref{thm:D} is the argument below, about two thousand lines of Python, and the checks
listed in Section~\ref{sec:checks}, among them two reports by fresh instances of the same model that were asked
to break it and found no error that affects the theorem. During the work one real gap was found, by a numerical
comparison and not by me: the pieces of Section~\ref{sec:cert} did not cover one family of pairs
(Remark~\ref{rem:gap}). The result rests on \cite{NY,N2} (formalised) and on the theorem of Coven and
Meyerowitz. I have not established that it is new: the paper that posed the question says of $M=p^\alpha q^\beta$
``we do not know''~\cite[Section~4]{KLMS}, and web searches made for \cite{N2} on 20 November 2026 found no proof
for any exponents beyond those of \cite{NY}; that is less than a specialist's knowledge of the field. Lemma Y$'$ of
\cite{N2}, whose mechanism reappears in Lemma~\ref{lem:M}, is due to Yonder, another instance on the same
machine, who also confirmed the theorem in five of the groups by a search over sets that uses nothing from this
paper.}}

\section{Introduction}

For $A\subseteq\Z_M$ a \emph{weak tiling complement} is a function $f:\Z_M\to\R$ with $f\ge0$, $f(0)=1$ and
$\one_A*f=1$, that is, $\sum_{a\in A}f(v-a)=1$ for every $v$; if one exists, $A$ is a \emph{weak tile}
\cite{LM,KLMS}. A tile is a weak tile, and a \emph{pd-tile}, a weak tile with a complement whose Fourier
transform is nonnegative, lies in between; spectral sets are pd-tiles \cite{KLMS}. Kiss, Londner, Matolcsi and
Somlai ask whether every pd-tile of $\Z_M$ is a tile and prove it for prime powers; Kadir and Fan~\cite{KF} prove
it for $M=pq$. Every \emph{weak} tile is a tile for $M=p^nq$ and $M=pqr$ \cite{NY} and for $M=p^2q^2$ \cite{N2}.

\begin{theorem}\label{thm:D}
Let $p\ne q$ be primes. Every weak tile of $\Z_{p^3q^2}$ satisfies the conditions \textup{(T1)} and \textup{(T2)}
of Coven and Meyerowitz and has no orphan zero. In particular it is a tile.
\end{theorem}

So in these groups tiles, pd-tiles and weak tiles are the same sets. The smallest of the groups, $\Z_{72}$, is the
smallest cyclic group in which a tiling need not have a periodic factor.

The proof for $p^2q^2$ in \cite{N2} sorted the possible sets $Z$ of cyclotomic divisors (``patterns'') of a weak
tile by a linear programme and needed one counting lemma for two patterns. Carried over as it stands to
$p^3q^2$ it leaves dozens of patterns undecided for every pair of primes. Two things are added here.

\emph{Reductions at the top} (Section~\ref{sec:top}). The reductions of \cite{N2} work with the subgroups
$p\Z_M$ and $q\Z_M$ of prime index. The same can be done with the subgroups $P$ and $Q$ of prime order: a weak
tile that is a union of cosets of $P$ comes from $\Z_{M/p}$, and so does one with a weak complement equal to $1$
on $P$. If $M$ itself is a zero of $A$, then $A$ meets every coset of $P+Q$ in a union of cosets of $P$ or of $Q$;
so a weak tile that escapes the reductions is \emph{line-built}: a union of cosets of $P$ and of cosets of $Q$.
A line-built set is two smaller sets glued together, and three counting lemmas (Lemma~\ref{lem:M}) exclude the
ones that occur.

\emph{Certificates in two parameters} (Section~\ref{sec:cert}). Which reduction applies to a pattern is decided
by the linear programme, whose data are polynomials in $p$ and $q$. For $p^2q^2$ two regimes sufficed. Here the
plane of pairs $(p,q)$ is cut into ten pieces, on each of which every pattern has a certificate with $p$ and $q$
as symbols: a basis whose primal and dual solutions are rational functions with certified signs. In $73$ cases
the optimal basis changes inside a piece, and the certificate is a small tree of case distinctions on the sign
of a polynomial.

Sections~\ref{sec:pre} to~\ref{sec:M} contain the mathematics and are meant to be checkable by hand.
Section~\ref{sec:sort} states exactly what the computer has to establish; Section~\ref{sec:cert} says how it
does; Section~\ref{sec:checks} says what was checked, by whom, and what was not.

\section{Notation, and what is used from \texorpdfstring{\cite{NY,N2}}{[NY], [N2]}}\label{sec:pre}

Let $M=p^aq^b$ with $a,b\ge1$, and $\Z_M\cong\Z_{p^a}\times\Z_{q^b}$, $x\mapsto(u,w)$. The \emph{class} of $x$ is
$(i,j)$ with $i=\min(v_p(u),a)$, $j=\min(v_q(w),b)$; the class $(a,b)$ is $\{0\}$ and $(0,0)$ is the set of
units. A function that is constant on classes is a \emph{step function}. We write the order $p^\ell q^m$ of a
character as $(\ell,m)$, and say that $(\ell,m)\ne(0,0)$ is a \emph{zero} of $A\subseteq\Z_M$ if $\hat\one_A$
vanishes at the characters of that order, equivalently if the cyclotomic polynomial $\Phi_{p^\ell q^m}$ divides
$A(X)=\sum_{a\in A}X^a$. Let $Z(A)$ be the set of zeros,
\[
 E_p=\{\ell\ge1:(\ell,0)\in Z(A)\},\qquad E_q=\{m\ge1:(0,m)\in Z(A)\},\qquad k=p^{|E_p|}q^{|E_q|},
\]
and $S_A$ the set of prime powers $p^\ell$, $\ell\in E_p$, and $q^m$, $m\in E_q$. An \emph{orphan} is a zero
$(\ell,m)$ with $\ell,m\ge1$, $\ell\notin E_p$ and $m\notin E_q$. The conditions of Coven and
Meyerowitz~\cite{CM} read: (T1) $|A|=k$; (T2) $(\ell,m)\in Z(A)$ whenever $\ell\in E_p$ and $m\in E_q$. For a
set $Z$ of orders we use the same words and write $k(Z)$.

For a step function with values $\gamma_{ij}$ the Fourier transform at a character of order $(\ell,m)$ is
\[
 \Lambda_{\ell m}(\gamma)=\sum_{i,j}\gamma_{ij}\,c_p(\ell,i)\,c_q(m,j),\qquad
 c_p(\ell,i)=\begin{cases}1&i=a,\\ p^{a-i}[\ell\le i]-p^{a-i-1}[\ell\le i+1]&i<a,\end{cases}
\]
and $c_q$ likewise; $\Lambda_{00}(\gamma)$ is the sum of the function.

\begin{itemize}
\item[(P1)] If $f$ is a weak complement of $A$, then $\hat f(\xi)\hat\one_A(\xi)=0$ for $\xi\ne0$ and
$|A|\sum f=M$. Conversely, if $h\ge0$, $h(0)=1$ and $\hat h(\xi)\hat\one_A(\xi)=0$ for all $\xi\ne0$, then
$\one_A*h$ is the constant $|A|\sum h/M$, which is at least $1$.
\item[(P2)] A weak tile has a weak complement that is a step function (a \emph{step complement}).
\item[(P3)] $k$ divides $|A|$. A set with (T1) and (T2) tiles $\Z_M$ \cite{CM}.
\item[(P4)] (\emph{Lemma 1.}) If $f$ is a weak complement of $A$ and $a\ne a'$ are in $A$, then $f(a-a')=0$.
\item[(P5)] (\emph{The bound} (S).) Let $g$ be a weak complement of $A$ and let $T\subseteq\Z_M$ be a set of
which every order $d>1$ with $\hat g\ne0$ at that order is a zero. Then $|T|\ge|A|$. (Indeed $\one_T*g$ is the
constant $|T|\sum g/M=|T|/|A|$, and at a point of $T$ it is at least $g(0)=1$.)
\item[(P6)] For a set $S$ of prime powers dividing $M$ the digit set $T_S$ has $\prod_{s\in S}p_s$ elements
($p_s$ the prime of $s$), and its zeros are the orders $d$ that are divisible exactly by some $s\in S$
\cite{CM}. So by (P5) with $T=T_{S_A}$, and (P3): \emph{a weak tile without orphan satisfies} (T1).
\end{itemize}

\begin{definition}
$\PP(M)$: every weak tile of $\Z_M$ satisfies (T1) and (T2) and has no orphan.
\end{definition}

$\PP(p^nq)$ holds for all $n$ and all primes $p\ne q$ \cite{NY} (``no orphan'' is the lower-bound lemma there
together with (T1)). $\PP(p^2q^2)$ is Lemma~\ref{lem:P} below.

\begin{proposition}[the linear programme]\label{prop:LP}
For a set $Z$ of orders let $LP(Z)$ be the programme in the unknowns $\gamma_{ij}\ge0$: $\gamma_{ab}=1$;
$\Lambda_{\ell m}(\gamma)=0$ for every order $(\ell,m)\notin Z$; minimise $\sigma(\gamma)=\Lambda_{00}(\gamma)$.
Let $\Dw(Z)=M/\min\sigma$. If $A$ is a weak tile with $Z(A)=Z$, then the step complements of $A$ are exactly the
optimal points of $LP(Z)$, and $|A|=\Dw(Z)$.
\end{proposition}

\begin{proof}
By (P1) a step complement is feasible with $\sigma=M/|A|$, and for a feasible point $\gamma$ the function
$\one_A*\gamma$ is the constant $|A|\sigma(\gamma)/M\ge1$, with equality exactly when $\gamma$ is a weak
complement.
\end{proof}

\section{Reductions at the bottom}\label{sec:bottom}

These are from \cite{N2}, stated for $M=p^aq^b$.

\begin{lemma}\label{lem:A}
If neither $p$ nor $q$ is a zero of a weak tile $A\ni0$ of $\Z_M$, then $A\subseteq p\Z_M$ or $A\subseteq q\Z_M$.
If $A\subseteq p\Z_M$, then $A'=\{y\in\Z_{M/p}:py\in A\}$ is a weak tile of $\Z_{M/p}$.
\end{lemma}

\begin{proof}
Let $f$ be a step complement, $\sigma=\sum f$. As $\hat\one_A$ does not vanish at the characters of order $p$,
$\hat f$ does, and the sum of $f$ over $p\Z_M$, the average of $\hat f$ over the characters of order dividing
$p$, is $\sigma/p$. Likewise the sum over $q\Z_M$ is $\sigma/q$. So the sum over $p\Z_M\cup q\Z_M$ is less than
$\sigma$, and $f>0$ at some unit, hence at all units. By (P4) no difference of two elements of $A$ is a unit. If
$a\in A\setminus p\Z_M$ and $b\in A\setminus q\Z_M$, then $q\mid a$ and $p\mid b$ (as $a-0$, $b-0$ are not
units), and $a-b$ is a unit. For the second statement, $f'(y)=f(py)$ is a weak complement of $A'$.
\end{proof}

\begin{lemma}\label{lem:B}
Let $A\subseteq p\Z_M$ and $A'$ as above. Then $(\ell,m)$ is a zero of $A$ if and only if either $\ell\ge1$ and
$(\ell-1,m)$ is a zero of $A'$, where $(0,m)$ is read as $q^m$ and $(0,0)$ is never a zero, or $\ell=0$ and
$(0,m)$ is a zero of $A'$. Consequently, if $A'$ satisfies \textup{(T1)}, \textup{(T2)} and has no orphan
relative to $M/p$, so does $A$ relative to $M$.
\end{lemma}

\begin{proof}
$A(X)=A'(X^p)$, and for a primitive $d$-th root $\zeta$ of unity $\zeta^p$ is primitive of order $d/p$ if
$p\mid d$ and of order $d$ otherwise. So $E_p(A)=E_p(A')+1$, $E_q(A)=E_q(A')$, a zero $(\ell,m)$ with
$\ell\ge2$, $m\ge1$ corresponds to the zero $(\ell-1,m)$ of $A'$, and $(1,m)$ is a zero of $A$ exactly when
$m\in E_q(A')$. (T1): $|A|=|A'|$. (T2) and ``no orphan'' follow by reading these equivalences in both directions.
\end{proof}

\begin{lemma}\label{lem:C}
Let $A\subseteq\Z_M$ and let its pieces $A^c=\{y\in\Z_{M/p}:c+py\in A\}$, $c\in\Z_p$, all satisfy
\textup{(T1)}, \textup{(T2)} and have no orphan relative to $M/p$, with the same set $S'$ of prime-power zeros.
Then $A$ satisfies \textup{(T1)}, \textup{(T2)} and has no orphan relative to $M$.
\end{lemma}

\begin{proof}
$A(X)=\sum_cX^cA^c(X^p)$. The pieces have equal size, so $p$ is a zero. If $p^2\mid d$ then
$\Phi_d(X)=\Phi_{d/p}(X^p)$, and comparing the parts of $A(X)$ in each residue class of exponents modulo $p$ one
sees that $d$ is a zero of $A$ if and only if $d/p$ is a zero of every piece. If $p\nmid m'$ and $m'$ is a zero
of every piece, then $m'$ and $pm'$ are zeros of $A$. Hence $S_A\supseteq\{p\}\cup\{p^{\ell+1}:p^\ell\in S'\}\cup
\{q^m\in S'\}$; the product of the primes of this set is $p\,|A^c|=|A|$, and $k$ divides $|A|$: equality, and
(T1). (T2) at $(\ell,m)$: for $\ell\ge2$ by (T2) of the pieces at $(\ell-1,m)$; for $\ell=1$ because $q^m$ is a
zero of every piece. No orphan: a zero $(1,m)$ is supported by $p$; a zero $(\ell,m)$ with $\ell\ge2$ gives the
zero $(\ell-1,m)$ of every piece, supported by $p^{\ell-1}$ or $q^m$ in $S'$.
\end{proof}

\begin{lemma}\label{lem:D}
Let $A_1,A_2$ be weak tiles of $\Z_{M'}$, $M'=p^{a'}q^{b'}$, with a common weak complement, both satisfying
\textup{(T1)} and without orphan. Then $S_{A_1}=S_{A_2}$.
\end{lemma}

\begin{proof}
Both have $n=M'/\sum g$ elements, and by (T1) $n=p^{\alpha_i}q^{\beta_i}$ with $\alpha_i,\beta_i$ the numbers of
powers of $p$ and of $q$ in $S_i=S_{A_i}$; so these numbers agree. Put $S_0=S_1\cap S_2$ and let $P_i$, $Q_i$ be
the powers of $p$, of $q$ in $S_i\setminus S_0$; $|P_1|=|P_2|=x$, $|Q_1|=|Q_2|=y$. Let $d>1$ be a zero of both
sets. Some $s_1\in S_1$ and $s_2\in S_2$ divide $d$ exactly. If neither lies in $S_0\cup P_1\cup P_2$, both are
powers of $q$ dividing $d$ exactly, hence equal, hence in $S_0$. So every common zero is a zero of the digit set
of $S_0\cup P_1\cup P_2$, and likewise of that of $S_0\cup Q_1\cup Q_2$. The transform of the common complement
vanishes wherever one of the two sets has no zero; so by (P5) and (P6) $n\le p^{2x}\prod_{S_0}p_s$ and
$n\le q^{2y}\prod_{S_0}p_s$, while $n=p^xq^y\prod_{S_0}p_s$. So $q^y\le p^x\le q^y$, and $x=y=0$.
\end{proof}

\begin{corollary}\label{cor:R}
Assume $\PP(M/p)$. If a weak tile $A$ of $\Z_M$ has a weak complement that vanishes outside $p\Z_M$, then $A$
satisfies \textup{(T1)}, \textup{(T2)} and has no orphan. The same with $q$.
\end{corollary}

\begin{proof}
With $g(y)=f(py)$, the sum $\sum_{a\in A}f(c+py-a)$ only sees the $a\equiv c$ modulo $p$; so every piece $A^c$
is a weak tile of $\Z_{M/p}$ with the weak complement $g$. Apply $\PP(M/p)$, Lemma~\ref{lem:D} and
Lemma~\ref{lem:C}.
\end{proof}

\begin{lemma}\label{lem:P}
$\PP(p^2q^2)$ holds.
\end{lemma}

\begin{proof}
Follow the proof of the theorem of \cite{N2} through its cases. If neither $p$ nor $q$ is a zero, Lemma~\ref{lem:A},
$\PP(pq^2)$, $\PP(p^2q)$ and Lemma~\ref{lem:B} give the three properties. If the pattern satisfies (T2) and has
no orphan, (T1) holds by (P6). In the cases where \cite{N2} shows $|A|<k$, or that $|A|$ is not an integer, or
applies the counting lemma, there is no weak tile. In the remaining cases a weak complement vanishes outside
$p\Z_M$ or $q\Z_M$, and Corollary~\ref{cor:R} applies.
\end{proof}

\section{Reductions at the top}\label{sec:top}

Let $P=(M/p)\Z_M$ and $Q=(M/q)\Z_M$ be the subgroups of order $p$ and $q$, and $K=P+Q$. For a step function
write $x,y,z$ for its values on $P\setminus\{0\}$, $Q\setminus\{0\}$ and $K\setminus(P\cup Q)$: the classes
$(a-1,b)$, $(a,b-1)$, $(a-1,b-1)$.

\begin{lemma}[lines]\label{lem:L}
Let $d=p^\ell q^m$ with $\ell,m\ge1$ be a zero of a multiset $A$ in $\Z_M$, and let $N(c)$ be the number of
elements of $A$ congruent to $c$ modulo $d$. On every coset of the subgroup of order $pq$ of $\Z_d$, written as
a $p\times q$ array, $N(u,v)=\alpha_u+\beta_v$ with integers $\alpha_u,\beta_v\ge0$. If $A$ is a set and $d=M$,
the array is a union of rows or a union of columns.
\end{lemma}

\begin{proof}
$\Phi_d$ divides $A(X)$ exactly when $(X^{d/p}-1)(X^{d/q}-1)A(X)$ vanishes modulo $X^d-1$, because the product
contains every cyclotomic factor of $X^d-1$ other than $\Phi_d$; that is, when
$N(c)-N(c+d/p)-N(c+d/q)+N(c+d/p+d/q)=0$ for all $c$. On a coset this gives
$N(u,v)=N(u,v_0)+N(u_0,v)-N(u_0,v_0)$; take $u_0$ where $N(\cdot,v_0)$ is least,
$\alpha_u=N(u,v_0)-N(u_0,v_0)$ and $\beta_v=N(u_0,v)$. An array of zeros and ones of this form has all
$\alpha_u=0$ or all $\beta_v=0$.
\end{proof}

This is the two-prime case of the theorem of Lam and Leung on vanishing sums of roots of unity~\cite{LL}.

\begin{lemma}\label{lem:E}
Let $A$ be a weak tile of $\Z_M$ and $f$ a step complement of $A$.
\begin{enumerate}
\item $f\le1$ everywhere.
\item $\Lambda_{ab}(f)=1-x-y+z$. So if $M$ is not a zero of $A$, then $x+y=1+z$.
\item If $x=1$, then $A$ meets every coset of $P$ at most once, and its image $A'$ in $\Z_M/P\cong\Z_{M/p}$ is
a weak tile. The zeros of $A$ among the orders $(\ell,m)$ with $\ell<a$ are exactly the zeros of $A'$. If $A'$
satisfies \textup{(T1)} and \textup{(T2)} relative to $M/p$, then $p^a$ is not a zero of $A$, and $A$ satisfies
\textup{(T1)} and \textup{(T2)} relative to $M$.
\item If $A+P=A$, the image $A'$ of $A$ in $\Z_{M/p}$ is a weak tile, and
$Z(A)=Z(A')\cup\{(a,m):0\le m\le b\}$. \textup{(T1)}, \textup{(T2)} and ``no orphan'' pass from $A'$ to $A$.
\item If $M$ is a zero of $A$, then $A$ meets every coset of $K$ in a union of cosets of $P$ or in a union of
cosets of $Q$.
\end{enumerate}
The same holds with $P$, $p$, $x$ exchanged with $Q$, $q$, $y$.
\end{lemma}

\begin{proof}
(1) For $a\in A$ and any $g$: $1=\sum_{a'\in A}f(a+g-a')\ge f(g)$. (2) A character of order $M$ is non-trivial
on $P$ and on $Q$; its sums over $P\setminus\{0\}$, $Q\setminus\{0\}$ and $K\setminus(P\cup Q)$ are $-1$, $-1$
and $1$, and it sums to zero over every other class except $\{0\}$. (3) For $a\in A$ and
$g\in P\setminus\{0\}$: $1=\sum_{a'}f(a+g-a')\ge f(g)=1$, so $f(a-a'+g)=0$ for all $a'\ne a$ and all $g\in P$;
with (P4), $a-a'\notin P$. Let $f'$ be the push-forward of $f$ to $\Z_M/P$ divided by $p$. Its value at $0$ is
$(1+(p-1)x)/p=1$, and at the image of $v$ the sum of $f'$ over the differences with the points of $A'$ is
$\frac1p\sum_{g\in P}(\one_A*f)(v+g)=1$. A character of order $(\ell,m)$ with $\ell<a$ is trivial on $P$, so it
takes the same values on $A$ and on $A'$. If $A'$ has (T1), then $|A|=|A'|$ is the product of the primes of the
prime-power zeros of $A'$; a further zero $p^a$ of $A$ would multiply the number dividing $|A|$ by $p$. So the
prime-power zeros of $A$ are those of $A'$, all with $\ell<a$, and (T1), (T2) for $A$ are the statements for
$A'$. (4) Here $x=0$ by (P4), so the push-forward $f'$ has $f'(0)=1$, and $\one_{A'}*f'$ at the image of $v$ is
$\sum_{a\in A}f(v-a)=1$. A character of order $(a,m)$ is non-trivial on $P$, so its sum over every coset of
$P$, and hence over $A$, is zero; the other orders are as in (3). $S_A=S_{A'}\cup\{p^a\}$ and $|A|=p\,|A'|$ give
(T1); (T2) at $(a,m)$ holds because every $(a,m)$ is a zero, and elsewhere it is inherited, like ``no orphan''.
(5) Lemma~\ref{lem:L} with $d=M$.
\end{proof}

In (3) ``no orphan'' is not obtained for $A$: a zero $(a,m)$ of $A$ whose $q^m$ is not a zero would be one. This
is why the use of (3) in Section~\ref{sec:sort} needs an extra condition on the pattern.

\begin{corollary}\label{cor:TU}
Assume $\PP(M/p)$ and $\PP(M/q)$, and let $A$ be a weak tile of $\Z_M$.
\begin{itemize}
\item[(T)] If some step complement of $A$ has $x=1$ or $y=1$, then $A$ satisfies \textup{(T1)} and \textup{(T2)}.
\item[(U)] If $M$ is a zero of $A$ and some step complement of $A$ has $x>0$ or $y>0$, then $A$ satisfies
\textup{(T1)}, \textup{(T2)} and has no orphan.
\item[(H)] If $M$ is a zero of $A$ and $Z(A)$ violates \textup{(T2)} or has an orphan, then $A$ is not a union
of cosets of $P$, nor of cosets of $Q$; so $A$ is the union of $i\ge1$ cosets of $P$ and $j\ge1$ cosets of $Q$
that lie in different cosets of $K$.
\end{itemize}
\end{corollary}

\begin{proof}
(T) is Lemma~\ref{lem:E}(3). (U): if $x>0$, no two elements of $A$ differ by an element of
$P\setminus\{0\}$ (P4), so $A$ contains no coset of $P$; by Lemma~\ref{lem:E}(5) it meets every coset of $K$ in
a union of cosets of $Q$, so $A+Q=A$, and Lemma~\ref{lem:E}(4) applies. (H): by Lemma~\ref{lem:E}(4) a union of
cosets of $P$ or of $Q$ would satisfy (T2) and have no orphan; the rest is Lemma~\ref{lem:E}(5), and two cosets
of different kinds in one coset of $K$ would intersect.
\end{proof}

\begin{remark}
If $M$ is not a zero of $A$ and no step complement has $x=1$ or $y=1$, then by (1) and (2) every step complement
has $x,y>0$, and $A$ meets every coset of $P$ and of $Q$ at most once. No pattern of $\Z_{p^3q^2}$ ends in this
case, as the computation shows; in $\Z_{p^4q^2}$ some do (Section~\ref{sec:rem}).
\end{remark}

\section{Line-built sets in \texorpdfstring{$\Z_{p^3q^2}$}{Z\_\{p\^{}3q\^{}2\}}}\label{sec:M}

From here on $M=p^3q^2$. Call $A\subseteq\Z_M$ \emph{line-built} if it is the union of $i\ge1$ cosets of $P$
and $j\ge1$ cosets of $Q$ that lie in different cosets of $K$. Let $X\subseteq\Z_M/P\cong\Z_{p^2q^2}$ and
$Y\subseteq\Z_M/Q\cong\Z_{p^3q}$ be the sets of these cosets. A character of order $(\ell,m)$ is trivial on $P$
exactly when $\ell\le2$ and on $Q$ exactly when $m\le1$, and its sum over a coset of a subgroup on which it is
not trivial vanishes. So for a character $\chi$ of order $(\ell,m)$:
\begin{equation}\label{eq:dict}
\chi(A)=\begin{cases}0&(\ell,m)=(3,2),\\ q\,\chi(Y)&\ell=3,\ m\le1,\\ p\,\chi(X)&\ell\le2,\ m=2,\\
p\,\chi(X)+q\,\chi(Y)&\ell\le2,\ m\le1,\end{cases}
\end{equation}
and, as a coset of $P$ or of $Q$ lies in one class modulo $p^2q$, the numbers of elements in a class modulo
$p^\ell q^m$ with $\ell\le2$, $m\le1$ satisfy $N^A=p\,N^X+q\,N^Y$.

\begin{lemma}\label{lem:M}
Let $A\subseteq\Z_{p^3q^2}$ be line-built.
\begin{itemize}
\item[(M0)] $|A|=ip+jq$. So $|A|\ne pq$; and if $|A|=p^2q$, then $j=pj'$ and $i=q(p-j')$ with $1\le j'\le p-1$.
\item[(M2)] If $|A|=p^2q$ and $p$, $p^3q$ and $p^2q^2$ are zeros of $A$, then $A+K=A$.
\item[(M4)] If $|A|=p^2q$, then $q$ and $pq^2$ are not both zeros of $A$.
\end{itemize}
\end{lemma}

\begin{proof}
(M0) If $ip+jq=pq$ then $p$ divides $j$, so $jq\ge pq$. If $ip+jq=p^2q$ then $p$ divides $j$, and $i\ge1$ gives
$j'<p$.

(M2) As $p$ is a zero, every class $c$ modulo $p$ holds $pq=p\,N^X(c)+q\,N^Y(c)$ elements; so $N^Y(c)=p\,t_c$
and $N^X(c)=q(1-t_c)$ with $t_c\in\{0,1\}$. By \eqref{eq:dict}, $p^3q$ is a zero of $Y$, the top order of
$\Z_{p^3q}$; so by Lemma~\ref{lem:L} the set $Y$ meets every coset of the subgroup of order $pq$ of $\Z_{p^3q}$
in a union of cosets of its subgroup of order $p$ or of its subgroup of order $q$. Such a coset lies in one
class modulo $p$. A class with $t_c=1$ holds $p$ points of $Y$, a sum of multiples of $p$ and multiples of $q$;
a positive number of cosets of order $q$ would have to be a multiple of $p$ and hold at least $pq$ points. So
these $p$ points are one coset of the subgroup of order $p$, which is $K/Q$: the cosets of $Q$ in $A$ fill whole
cosets of $K$. In the same way $p^2q^2$ is a zero of $X$, the top order of $\Z_{p^2q^2}$, and a class with
$t_c=0$ holds $q$ points of $X$: one coset of the subgroup of order $q$, which is $K/P$. So $A+K=A$.

(M4) As $q$ is a zero, every class $t$ modulo $q$ holds $p^2=p\,N^X(t)+q\,N^Y(t)$ elements; so $N^Y(t)=p\,s_t$
and $N^X(t)=p-q\,s_t\ge0$. Let $pq^2$ be a zero of $A$, hence of $X$ by \eqref{eq:dict}. The class $t$ modulo
$q$ is a coset of the subgroup of order $pq$ of $\Z_{pq^2}$, and by Lemma~\ref{lem:L} the numbers of elements of
$X$ in the classes modulo $pq^2$ inside it are a sum of $\alpha$ full rows, of $q$ points, and $\beta$ full
columns, of $p$ points: $p-q\,s_t=\alpha q+\beta p$. If $\beta=0$, $q$ divides $p$. So $\beta=1$, $\alpha=0$ and
$s_t=0$, for every $t$; and $Y$ is empty.
\end{proof}

If $A+K=A$ then $A$ is a union of cosets of $P$. The two patterns for which \cite{N2} needed Yonder's counting
lemma are of the kind (M0): there $M=p^2q^2$ is a zero and $|A|=pq$.

\section{What has to be computed}\label{sec:sort}

Let $M=p^3q^2$. There are $2^{11}=2048$ sets $Z$ of orders. In the classes of Section~\ref{sec:top},
$P\setminus\{0\}$ is the class $(2,2)$ and $Q\setminus\{0\}$ is $(3,1)$.

\begin{definition}\label{def:cat}
A set $Z$ of orders of $\Z_{p^3q^2}$ is of category
\begin{itemize}
\item[\textbf{L2}] if neither $(1,0)$ nor $(0,1)$ is in $Z$;
\item[\textbf{a}] if $Z$ satisfies (T2) and has no orphan (and is not of category L2).
\end{itemize}
Otherwise it is \emph{settled} at the pair $(p,q)$ if one of the following holds.
\begin{itemize}
\item[\textbf{b}] $\Dw(Z)<k(Z)$.
\item[\textbf{N}] $\Dw(Z)/k(Z)$ is not an integer.
\item[\textbf{R}] $LP(Z)$ has an optimal point that vanishes on the classes $(0,j)$, or one that vanishes on
the classes $(i,0)$.
\item[\textbf{T}] $(3,2)\notin Z$, $Z$ violates (T2), and $LP(Z)$ has an optimal point with $\gamma_{22}=1$ or
with $\gamma_{31}=1$.
\item[\textbf{U}] $(3,2)\in Z$ and $LP(Z)$ has an optimal point with $\gamma_{22}>0$ or with $\gamma_{31}>0$.
\item[\textbf{H}] $(3,2)\in Z$, and either $\Dw(Z)=pq$; or $\Dw(Z)=p^2q$ and $(1,0),(3,1),(2,2)\in Z$; or
$\Dw(Z)=p^2q$ and $(0,1),(1,2)\in Z$.
\end{itemize}
\end{definition}

There are $512$ sets of category L2 and $321$ of category a, whatever the primes; $1215$ remain.

\begin{proposition}\label{prop:sort}
Let $p\ne q$ be primes such that each of the $1215$ sets $Z$ outside the categories \textup{L2} and \textup{a}
is settled at $(p,q)$. Then $\PP(p^3q^2)$ holds for these primes; moreover no weak tile of $\Z_{p^3q^2}$ has a
set of zeros outside \textup{L2} and \textup{a}.
\end{proposition}

\begin{proof}
Let $A$ be a weak tile and $Z=Z(A)$; recall $|A|=\Dw(Z)$ (Proposition~\ref{prop:LP}), that $k(Z)$ divides $|A|$,
and that $\PP$ holds for $\Z_{p^2q^2}$ (Lemma~\ref{lem:P}) and $\Z_{p^3q}$.
L2: by Lemmas~\ref{lem:A} and~\ref{lem:B}, $A$ has the three properties. a: (T2) and ``no orphan'' hold by
definition and (T1) by (P6). Otherwise $Z$ violates (T2) or has an orphan, and we show that this is impossible.
b, N: $k$ divides $|A|$. R: the optimal point is a weak complement of $A$ (Proposition~\ref{prop:LP}) vanishing
outside $p\Z_M$ or outside $q\Z_M$; by Corollary~\ref{cor:R}, $A$ satisfies (T2) and has no orphan. T: by
Corollary~\ref{cor:TU}(T), $A$ satisfies (T2). U: by Corollary~\ref{cor:TU}(U), $A$ satisfies (T2) and has no
orphan. H: by Corollary~\ref{cor:TU}(H), $A$ is line-built, with $|A|=\Dw(Z)$; (M0) excludes $|A|=pq$; under
the second condition (M2) gives $A+K=A$, which Corollary~\ref{cor:TU}(H) has excluded; under the third (M4)
applies.
\end{proof}

So Theorem~\ref{thm:D} follows once every set $Z$ outside L2 and a is settled at every pair of distinct primes.

\section{The computation}\label{sec:cert}

All files are in the directory \file{p2q2-proof-agent2/next-case-p3q2/} of the collection this paper belongs to;
\file{param/README.md} there lists them, and the working note \file{p2q2-proof-agent2/towards-paqb-agent2.md}
records the order in which things were found, mistakes included.

\subsection{Single pairs} For a given pair of primes, \file{exact.py} sorts the $2048$ sets, solving every
linear programme by the simplex method in rational arithmetic, and \file{settle.py} checks the conditions of
category H. This was done for all $72$ ordered pairs of distinct primes up to $23$ and for $27$ further pairs;
every set is settled at each of them (\file{settle72.out}, \file{settle-below.out}). For Theorem~\ref{thm:D} six
of these pairs are needed.

\subsection{Certificates with symbols}\label{ss:sym} The coefficients of $LP(Z)$ are polynomials in $p$ and
$q$. A \emph{certificate} for $Z$ on a set $\Omega$ of real pairs $(p,q)$ consists of
\begin{itemize}
\item a \emph{basis}: a choice of unknowns, as many as there are independent equations, such that the basic
solution $\gamma(p,q)$, a vector of rational functions, is $\ge0$ on $\Omega$, and such that the dual basic
solution $y(p,q)$ has reduced costs $\ge0$ on $\Omega$, neither $\gamma$ nor $y$ having a pole on $\Omega$. By
linear-programming duality $\sigma(\gamma)$ is then the optimal value at every point of $\Omega$, so $\Dw(Z)$ is
a known rational function there;
\item and for the category: (b) $k-\Dw>0$ on $\Omega$; (R, T, U) a second basis of the programme with the
extra requirement (some unknowns removed; or the equation $\gamma_{22}=1$ or $\gamma_{31}=1$ added; or the
equation $\sigma=M/\Dw$ added) whose basic solution is $\ge0$ on $\Omega$, has the optimal value, and has the
property required, for U strictly; (H) $\Dw=pq$ or $\Dw=p^2q$ as rational functions, and the conditions on $Z$;
(N) an argument that $\Dw/k$ is not an integer at any pair of distinct primes in $\Omega$, of one of the forms
below.
\end{itemize}
The signs are certified as follows. $\Omega$ is given by a parametrisation $(p,q)=\phi(r,s)$ or $\phi(r,z)$ with
parameters $\ge0$ (see the table below). A polynomial $g$ is \emph{positive on $\Omega$} if, after substituting
$\phi$ and multiplying by the power of $1+z$ that clears the denominator, all its coefficients are $\ge0$ and its
value at $r=s=z=0$ is $>0$. A rational function is signed through the irreducible factors of its numerator and
denominator.

\subsection{Integrality}\label{ss:int} Let $r=\Dw/k$, a rational function with integer coefficients in lowest
terms. That $r(p,q)$ is not an integer for primes $p\ne q$ in $\Omega$ is certified in one of three ways.
\begin{itemize}
\item[(iii)] $c<r<c+1$ on $\Omega$ for an integer $c$.
\item[(ii)] One of the primes, say $q$, divides the denominator as a polynomial, and the numerator at $q=0$ is
$\pm p^e$: an integer value would give $q\mid p^e$. (No certificate in the end uses this form.)
\item[(i)] Write $r=p^iq^jG/H$ with $G,H$ polynomials not divisible by $p$ or $q$. For each of the two primes
$x$, with $y$ the other, $H\equiv\pm y^eu(y)$ modulo $x$ for a polynomial $u$ with positive leading
coefficient, and one of: $u=1$; or $mx<u(y)<(m+1)x$ on $\Omega$ for an integer $m\ge0$; or $mx<u(y)<(m+2)x$ on
$\Omega$ and $u$ is linear. In the first two cases $x\nmid H(p,q)$. In the third, $x\mid H(p,q)$ forces
$u(y)=(m+1)x$; on that curve $r$ is a rational function $r_1$ of $x$ alone, $r_1=Q_1+R_1$ with $Q_1$ a polynomial
and $R_1\to0$, and with $c$ a common denominator of the coefficients of $Q_1$ one has $0<|c\,R_1(x)|<1$ for
$x\ge T_1$ (decided by isolating real roots), so $r_1(x)$ is not an integer there; for $2\le x<T_1$ it is
evaluated. Off the curves $H(p,q)$ is prime to $p$ and $q$, so it would divide $G(p,q)$; and
$0<|G-cH|<|H|$ on $\Omega$ for an integer $c$ excludes that.
\end{itemize}
Of the $164$ certificates of category N, $44$ use (iii) and $120$ use (i); among the latter three need an
inequality $mx<u(y)<(m+1)x$ and one a curve ($q=2p-1$, where $r_1=q(q+3)/(3q+1)$); in the others $u=1$.

\subsection{Case distinctions}\label{ss:split} If no basis is certified on the whole of a piece $\Omega$, the
generator takes a polynomial factor $F$ whose sign it could not certify and treats
$\Omega\cap\{F\ge0\}$ and $\Omega\cap\{F\le0\}$ separately, and again inside them (to depth two at most). On
such a part $F$ counts as weakly signed, and a weakly signed factor may not occur in a denominator. A part in
which no prime pair was found is closed by a proof that it is empty: polynomials $m_i$ with coefficients
$\ge0$ such that $\sum_im_is_iF_i$, where $s_iF_i\ge0$ are the assumptions, is negative on $\Omega$. (This
happens once: on the piece A1b, for the two factors $F_1$, $F_2$ of one tree, $q(-F_1)+F_2=p^2(q-1)(p-q)<0$.)

\subsection{The pieces}\label{ss:pieces}
\begin{center}\small
\begin{tabular}{llrrrr}
\toprule
piece & range & one certificate & trees & leaves & empty \\
\midrule
A1a & $p\ge5$, $p+2\le q<2p-\tfrac52$ & 1210 & 5 & 1220 & 0\\
A1b & $p\ge5$, $2p-1\le q<p^2-p+1$ & 1156 & 59 & 1275 & 1\\
A2 & $p\ge3$, $p^2-p+1\le q<p^2-1$ & 1214 & 1 & 1216 & 0\\
B1 & $p\ge5$, $p^2+1\le q<p^2+p$ & 1215 & 0 & 1215 & 0\\
B2 & $p\ge3$, $p^2+p+1\le q<p^3-p^2+p-1$ & 1213 & 2 & 1217 & 0\\
B3 & $p\ge3$, $p^3-p^2+p\le q<p^3-1$ & 1214 & 1 & 1216 & 0\\
C$_3$ & $p\ge3$, $q\ge p^3+1$ & 1215 & 0 & 1215 & 0\\
Ap1 & $q\ge3$, $q+1\le p<q^2-q+1$ & 1210 & 5 & 1220 & 0\\
Ap2 & $q\ge3$, $q^2-q+1\le p<q^2-1$ & 1215 & 0 & 1215 & 0\\
C$'$ & $q\ge2$, $p\ge q^2+1$ & 1215 & 0 & 1215 & 0\\
\bottomrule
\end{tabular}
\end{center}
A band $\mathit{lo}\le t<\mathit{up}$ is parametrised by $t=(\mathit{lo}+\mathit{up}\cdot z)/(1+z)$, $z\ge0$,
which reaches $\mathit{lo}$ and never $\mathit{up}$; the other prime is $5+r$, $3+r$ or $2+r$; C$_3$ is
$q=p^3+1+s$ and C$'$ is $p=q^2+1+s$. In addition, for $p=2$ the certificates were computed with $q=t$ as the
only symbol and hold for all real $t\ge9$ (\file{family.py}; the sign of each rational function of $t$ beyond a
threshold is decided by isolating real roots).

\begin{lemma}[coverage]\label{lem:cover}
Every pair of distinct primes other than $(2,3)$, $(2,5)$, $(2,7)$, $(3,2)$, $(3,5)$, $(3,11)$ lies in one of
the ten pieces or has $p=2$, $q\ge11$.
\end{lemma}

\begin{proof}
Let $p\ge5$ and $q>p$. Then $q=p+1$ and $q=2p-2$ are even; the integers from $p+2$ to $2p-3$ lie in A1a, those
from $2p-1$ to $p^2-p$ in A1b, those from $p^2-p+1$ to $p^2-2$ in A2; $p^2-1$ and $p^2$ are composite; the
integers from $p^2+1$ to $p^2+p-1$ lie in B1; $p^2+p$ is composite; those from $p^2+p+1$ to $p^3-p^2+p-2$ lie in
B2; $p^3-p^2+p-1=(p-1)(p^2+1)$ is composite; those from $p^3-p^2+p$ to $p^3-2$ lie in B3; $p^3-1$ and $p^3$ are
composite; and $q\ge p^3+1$ is C$_3$. For $p=3$ the same holds from A2 on ($q=7$; $13,17,19$; $23$; $q\ge29$),
which leaves $q=5,11$. For $q\ge3$ and $p>q$: the integers from $q+1$ to $q^2-q$ lie in Ap1, those from
$q^2-q+1$ to $q^2-2$ in Ap2, $q^2-1$ and $q^2$ are composite, and $p\ge q^2+1$ is C$'$. For $q=2$: $p=3$ is
excepted and $p\ge5$ is in C$'$. For $p=2$: $q=3,5,7$ are excepted.
\end{proof}

\begin{remark}\label{rem:gap}
In a first version the band A1a ended at $2p-3$, and the numerical comparison of Section~\ref{sec:checks}
treated bands as closed. It reported two mismatches at $(p,q)=(7,43)$, the upper end $q=p^2-p+1$ of A1b, which
the parametrisation does not reach and which belongs to A2. Looking at the ends again showed that nothing
covered $q=2p-3$. The bands are half-open, and Lemma~\ref{lem:cover} was written after this.
\end{remark}

\subsection{Proof of Theorem \ref{thm:D}} By Proposition~\ref{prop:sort} it suffices that every set $Z$ outside
L2 and a is settled at every pair of distinct primes. At the six excepted pairs this is the computation for
single pairs; for $p=2$, $q\ge11$ the certificates in $t$; for all other pairs, by Lemma~\ref{lem:cover}, the
certificates on the piece that contains the pair, where for a tree the pair lies in at least one part. In every
certificate of category T the set $Z$ violates (T2), and in those of categories U and H the order $(3,2)$ is in
$Z$, in those of T it is not; this is checked with the certificates. \qed

\section{What was checked, and what was not}\label{sec:checks}

\emph{By programs.} (1) \file{verify3.py} and \file{verify\_fast.py} import nothing from the generators. For every
set $Z$ on every piece they recompute both basic solutions from the stored bases, check that they solve their
equations, that no component has a pole on the piece, and every sign condition, with their own implementation of
the sign test; they re-establish the integrality arguments, walk the trees and re-check the emptiness proof;
and they check that all $2048$ sets are present and that those of categories T, U, H are as the proof needs.
Result: ten pieces, $12\,224$ leaves, no failure; the certificates in $t$ for $p=2$: $1215$, no failure. (2)
\file{crosscheck3.py} and \file{crosscheck2.py} compare, at $139$ prime pairs in pieces, the certified optimum
with the optimum found by the exact simplex of \file{exact.py}, and test the property of the category there:
$168\,885$ comparisons, no mismatch after the correction of Remark~\ref{rem:gap}. (3) The sorting of
\file{exact.py} was also computed in floating point with a different solver, with the same counts.

\emph{By the first referee} (a fresh instance; report in \file{referee/report.md}; it saw the lemmas and the
computation for $34$ single pairs, before the certificates with symbols existed). It re-derived
Lemmas~\ref{lem:B} to~\ref{lem:M} and Proposition~\ref{prop:sort} by hand, recomputed the sorting at the $34$
pairs with its own exact simplex and a primal and dual certificate for every optimum ($69\,632$ sets, no
difference), and checked Lemma~\ref{lem:M} by listing all line-built sets of the relevant sizes in $\Z_{72}$ and
$\Z_{108}$ ($1\,137\,312$ and $8\,537\,991\,192$ sets). It found no error affecting the result, a harmless slip
in my simplex, and that $\PP(p^2q^2)$ was used without a written proof (now Lemma~\ref{lem:P}).

\emph{By the second referee} (another fresh instance; \file{referee2/report.md}; it was asked to attack the
certificates and the reasoning that generator and checker share). Its own exact simplex at $933$ prime pairs of
its choosing, among them every ordered pair of primes below $45$, the ends of every band and both sides of
every branching curve: $1\,126\,305$ certificate checks, no mismatch. Lemma~\ref{lem:cover} by brute force for
the $184\,470$ ordered pairs of distinct primes below $3000$. Every certificate evaluated in integer arithmetic
at $183\,088$ prime pairs. The integrality arguments re-derived by hand. It found no error affecting the
theorem; it found that my checkers did not look at the poles of the dual solution, did not check the solutions
against their equations, and did not themselves verify one condition that the generator printed; it made these
checks, and they are now in the checkers of (1).

\emph{By Yonder.} A search over sets with a constraint solver, using none of the lemmas or categories of this
paper, finds no weak tile violating (T1) or (T2) in $\Z_{72}$, $\Z_{108}$, $\Z_{200}$, $\Z_{392}$ and $\Z_{500}$,
five patterns at the last two being settled by hand arguments, which I read and found correct.

\emph{Not checked.} No human has read this paper. Nothing here is formalised; the results it rests on
\cite{NY,N2} are. Generator and checkers both rely on the computer algebra system sympy for factoring
polynomials, for arithmetic in the field $\Q(p,q)$ and for isolating real roots; the numerical comparisons do
not depend on the first two. The second referee did not re-read the lemmas, and the first did not see the
certificates. A mistake in the reasoning that generator and checker share passes both; the one found
(Remark~\ref{rem:gap}) was found by a numerical comparison. Whether the theorem is new has not been
established.

\section{Remarks}\label{sec:rem}

1. \emph{The next level.} For $\Z_{p^4q^2}$ and $\Z_{p^3q^3}$ the same sorting, in floating point at a few pairs of
primes, leaves between $54$ and $216$ sets unsettled. Most of them again have $M$ as a zero, so that a weak tile
would be line-built and lemmas like (M0)--(M4) might apply. But in $\Z_{144}$ three have $M$ not a zero, the case
of the remark after Corollary~\ref{cor:TU}; no lemma here touches it, and it is the range ($p^4q^2$, $p<q<p^2$)
where, as reported in the literature (I have it from a summary of \cite{KLMS} and later work, not from
the papers), functions violate (T2).

2. \emph{Two folds.} If $p$ is a zero of $A$ and the zeros of $A$ are closed under $(0,m)\mapsto(1,m)$ where the
transform of a step complement $f$ does not vanish, then moving the mass of $f$ from the classes with $i=0$ onto
$p\Z_M$ gives a weak complement supported in $p\Z_M$; and there is a mirror image at the top, which makes $f$
constant on $P$. With either hypothesis proved for weak tiles, the linear programmes would not be needed. Both
are true a posteriori for the groups treated here and in \cite{NY,N2}; I have no proof of either that does not
go through the theorem.

3. \emph{Sizes.} The method needs $|A|=\Dw(Z)$ to be a specific number, $pq$ or $p^2q$, before the counting
lemmas apply. A direct argument that a line-built weak tile satisfies (T2), without the programme, would be the
natural replacement.

\begin{thebibliography}{9}
\bibitem{CM} E.~Coven and A.~Meyerowitz, \emph{Tiling the integers with translates of one finite set},
J.~Algebra 212 (1999), 161--174.
\bibitem{KF} M.~Kadir and K.~Fan, \emph{Tiles and weak tiles in $\Z_{pq}$}, arXiv:2607.02149 (2026).
\bibitem{KLMS} G.~Kiss, I.~Londner, M.~Matolcsi and G.~Somlai, \emph{Functional tilings and the
Coven--Meyerowitz tiling conditions}, arXiv:2411.03854; Discrete Analysis 2025:18.
\bibitem{LL} T.~Y.~Lam and K.~H.~Leung, \emph{On vanishing sums of roots of unity}, J.~Algebra 224 (2000),
91--109. (Cited from memory; not consulted. Lemma~\ref{lem:L} is proved here.)
\bibitem{LM} N.~Lev and M.~Matolcsi, \emph{The Fuglede conjecture for convex domains is true in all
dimensions}, Acta Math. 228 (2022), 385--420.
\bibitem{NY} Noon and Yonder, \emph{Weak tiles of $\Z_{p^nq}$ and $\Z_{pqr}$ are tiles: two short proofs, checked
in Lean}, 20 November 2026 (\file{paper/weak-tiles.pdf} in this folder; since extended by a third part).
\bibitem{N2} Noon, \emph{Weak tiles of $\Z_{p^2q^2}$ are tiles: one reduction and 192 linear programmes, checked
in Lean}, 20 November 2026 (\file{paper/weak-p2q2.pdf} in this folder).
\end{thebibliography}

\end{document}
